\section{Related Work}
\label{sec:related}

The study of LLM-based agents has progressed from role-specialized collaborations to increasingly modular systems oriented toward general assistance. Early frameworks such as CAMEL~\cite{li2023camel} and MetaGPT~\cite{hong2024metagpt} showed that assigning distinct roles and procedures elicits more disciplined reasoning, especially in structured domains like software engineering. As evaluations moved toward open-ended tasks, recent efforts emphasized broader empirical analysis and unified infrastructures: OAgents~\cite{oagent} systematically studies design choices for effective agents at scale, while AWorld~\cite{aworld2025} provides a unified playground that supports both computer- and phone-use tasks, encouraging iteration on the full stack of planning, tools, and evaluation. In parallel, Alita~\cite{qiu2025alita} advances the idea that tools themselves can be dynamically constructed and composed, reducing reliance on heavy predefinition and enabling self-evolution of agent capabilities. Complementary platforms (e.g., AutoAgent~\cite{tang2025autoagent}) lower the barrier to assembling agents and workflows, making orchestration more accessible without sacrificing modularity.

A second thread concerns the substrate that turns models into reliable systems: tools, retrieval, and memory. Toolformer~\cite{schick2023toolformer} demonstrated that models can self-learn to call APIs, and surveys such as Qin et al.\cite{qin2024tool} underscored that robust tool interfaces and execution traces are core to stability. Hierarchical research agents (e.g., Open Deep Research\cite{opendeepresearch} and DeepResearchAgent~\cite{DeepResearchAgent}) pair retrieval-centric planning with iterative decomposition, while memory mechanisms progress from reflective feedback~\cite{shinn2023reflexionlanguageagentsverbal} to hierarchical designs like A-Mem~\cite{xu2025mem} that separate working, semantic, and procedural layers—supporting long-horizon continuity and adaptive control. These system utilities collectively motivate architectures where communication is structured, tool calls are auditable, and prior experience can be retrieved and reused.

Finally, generalization across domains has become a central goal. OWL~\cite{owl2025} tackles this by decoupling domain-agnostic planning from domain-specific execution (WORKFORCE), training only the planner—via supervised trajectories and reinforcement learning—to transfer across new environments without retraining worker agents. Orthogonal efforts (e.g., AgentRefine~\cite{fu2025agentrefine}, TapeAgents~\cite{bahdanau2024tapeagentsholisticframeworkagent}, and MiroFlow~\cite{miroflow}) study reproducibility, refinement, and stability at scale. Benchmarks such as GAIA~\cite{mialon2023gaia} and BrowseComp~\cite{wei2025browsecomp} crystallize these trends by testing multimodal reasoning, browsing, and execution, revealing that while proprietary frameworks (e.g., Deep Research~\cite{deepresearch}, h2oGPTe~\cite{h2oGPTe2024h2oGPTe}) are strong, open-source systems are rapidly closing the gap. Distinct from the above, our work contributes a system-level design that intentionally fuses complementary agentic paradigms with hierarchical memory and a statistically validated tool suite, aiming for robust gains under real-world constraints.